\begin{claim}{D1.2}{Proved}{Vertical and logarithmic tangent bundles}
There is a canonical representable proper flat holomorphic map
$\bar f:X\to\mathcal B$ with $f=\rho\bar f$. If
\[
 \mathcal T=T_X(-\log W),\qquad
 \mathcal V=\ker\bigl(\mathcal T\to
                   \bar f^*T_{\mathcal B}(-p_0)\bigr),
\]
then, for rank-two vector bundles $E,G$ on $\mathcal B$,
\[
 \mathcal V=\bar f^*E,\quad E=\OO\oplus\OO(-1),\qquad
 G=R^1\bar f_*\OO_X=\OO\oplus\OO(-1),                 \tag{D.2}
\]
\[
 R^0\bar f_*\OO_X=\OO,\qquad R^2\bar f_*\OO_X=\OO(-1),
\]
and the following is an exact sequence of locally free sheaves on $X$:
\[
 0\longrightarrow\bar f^*E\longrightarrow\mathcal T
 \longrightarrow\bar f^*\OO(-5)\longrightarrow0.       \tag{D.3}
\]
The trivial summands in $E,G$ are generated by the vertical field $\xi$
and the constant integral cocycle $\gamma$, respectively.
\end{claim}
\begin{proof}
The divisors $S_j$ and their canonical sections give root data
$\OO_X(S_j)^{m_j}=f^*\OO_B(p_j)$, hence the indicated map to the root
curve. Over its order-$m_j$ chart the fibre product is the original
smooth torus family $J'_j$, since its free affine cyclic quotient is
exactly $f^{-1}(D_j)$. The comparison $J'_j\to X$ is finite etale;
in particular pullback is faithful and the stated local comparisons
descend. Away from the two root points the map is the original $f$.
Properness and flatness may be checked after this atlas: properness is
the properness of the original fillings, and flatness holds because a
nonconstant base parameter is a non-zero-divisor in every smooth
total-space local ring. Fibres have dimension two, including the cusp.

On a smooth torus chart a vertical field is a holomorphic linear
combination of the translation fields. Its coefficients are constant
on each compact fibre and therefore depend only on the base. The
vertical derivative of an affine transition is its linear part $R_g$;
the translating part, including the logarithmic regluing, has vertical
derivative the identity. At the cusp the two fields
$x_1\partial_{x_1},x_2\partial_{x_2}$ extend as logarithmic vector
fields in every toric chart. The deck maps multiply each $x_i$ by a
base-dependent unit and a power of the base parameter. Their vertical
logarithmic derivative fixes both coefficient directions. The
exponential overlap sends $\partial_{\zeta_i}$ to
$2\pi i x_i\partial_{x_i}$. These transitions depend only on the base
and define $E$; they prove the actual bundle identity
$\mathcal V=\bar f^*E$.

We check the kernel and surjectivity at every type of cusp point. At
$t=xy$, the kernel of $df$ is freely generated by
$x\partial_x-y\partial_y$ and $\partial_z$. At $t=xyz$ it is generated
by $x\partial_x-y\partial_y$ and $x\partial_x-z\partial_z$.
Indeed the equation $yz a+xz b+xy c=0$ in the regular local ring forces
$x\mid a,y\mid b,z\mid c$; after division the only relation is that
the three coefficients sum to zero. The logarithmic map sends
$\sum a_i z_i\partial_{z_i}$ to $(\sum a_i)t\partial_t$ and is
surjective. At a finite root point it is the ordinary relative tangent
sequence of the smooth map to $s$, not the singular differential of
$s^{m_j}$. These comparisons prove (D.3), since
$T_{\mathcal B}(-p_0)=\OO(7-12)$.

The invariant vector $e_2$ defines a subbundle $\OO\xi\subset E$.
It remains nonzero as a vector of the logarithmic vertical bundle at
the cusp, even where its image in the ordinary tangent fibre is zero.
The quotient line has automorphy factors $-1/\tau$ and $1/\tau$ under
the order-three and order-four generators. The modular discriminant
has weight twelve and no zero in the upper half-plane. Therefore its
reciprocal is a meromorphic section of the twelfth tensor power of
this quotient line. On the cusp chart
$\tau=(\log q)/(2\pi i)+h(q)$ with $h$ holomorphic, so the product
expansion of the discriminant is $q$ times a holomorphic unit.
There are no other zeros or poles, including on the finite root
charts. The quotient line thus has twelfth power
$\OO(-p_0)=\OO(-12)$, and the torsion-free Picard group in D1.1 makes
it $\OO(-1)$. Its extension by $\OO$ splits because
$H^1(\mathcal B,\OO(1))=0$.

For completeness the structure-sheaf computation at the cusp uses the
normalization $P=\mathrm dP_6$, the three double curves and the two
triple points. Their exact normalization complex has global section
complex
\[
 \C\xrightarrow{0}\C^3
 \xrightarrow{(a,b,c)\mapsto(a-b+c,a-b+c)}\C^2.
\]
Exactness at the two local models is the elementary gluing of functions
on the coordinate branches; all higher cohomology of the rational
normalization, the rational double curves and the points vanishes.
The complex therefore gives $(h^0,h^1,h^2)=(1,2,1)$. Its product is
the Cech product of the two-triangle torus and gives a nondegenerate
alternating pairing $H^1\wedge H^1\to H^2$. This also follows directly
from the two opposite oriented triangles: the two coordinate cocycles
have determinant one on their oriented fundamental cycle. These are
the local comparison and product in Claude and Alpoge (2026,
Lemma 9.12), used here on the same normalization and with both triple
points retained.

On all other geometric fibres of $\bar f$ the cohomology is that of a
smooth two-torus. Constancy of these three dimensions and proper flat
analytic base change give locally free direct images on the atlas.
Their canonical comparison maps commute with the finite group actions
and with overlap restriction; hence they descend to $\mathcal B$.
The determinant of (D.3), using $W=\bar f^*p_0$, gives
\[
 \omega_{X/\mathcal B}=\bar f^*(\det E)^{-1}.
\]
Relative Serre duality, applied on that same atlas and descended by
its trace pairing, gives $R^2\bar f_*\OO_X=\det E=\OO(-1)$.
Cup product is an isomorphism $\det G\to R^2\bar f_*\OO_X$, since it
is so on every geometric fibre by the preceding calculation.

The class $\gamma$ is nowhere zero in $G$. On a smooth marked fibre,
$H^1(\OO)$ is the quotient of constant lattice cocycles by the two
period rows. If $\gamma$ were a combination of those rows, evaluation
on $\widehat w,\widehat\delta$ would make both coefficients zero,
contradicting $\gamma(\widehat\gamma)=1$. The same argument applies
on each finite root cover, and affine translations act trivially on
these classes. At the cusp $\gamma$ vanishes on the angular sublattice
and descends to a nonzero character of its rank-two deck lattice.
In the normalization complex it is the corresponding nonzero
one-cocycle on the two-triangle torus. Its cup product with a complementary
integral character is nonzero, so it remains nonzero there. Thus
$\OO\gamma\subset G$ is a subbundle; its quotient is $\OO(-1)$ by
the determinant just computed. Once again $H^1(\OO(1))=0$ splits the
extension. This proves every assertion in (D.2).
\end{proof}
